\documentclass[10pt,a4paper]{amsart}
\usepackage[margin=19mm]{geometry}
\usepackage{amsmath,amssymb,mathtools}
\usepackage[scr=rsfs,cal=euler]{mathalfa}
\usepackage{xfrac}
\usepackage[unicode,hidelinks,bookmarks=false]{hyperref}
\newcommand{\ie}{i.e.\ }
\newcommand{\cf}{cf.\ }
\newcommand{\resp}{respectively\ }
\newcommand{\Subsection}{\subsection}
\providecommand{\nobreakdash}{\nobreak\hbox{-}}
\setlength{\emergencystretch}{2em}
\multlinegap0pt
\usepackage{bm}
\usepackage{bbm}
\usepackage{dsfont}
\usepackage{xspace}
\usepackage{cancel}
\usepackage{mathtools}
\usepackage{amsthm}
\makeatletter
\@ifundefined{definition}{\newtheorem{definition}{Definition}}{}
\@ifundefined{proposition}{\newtheorem{proposition}[definition]{Proposition}}{}
\@ifundefined{theorem}{\newtheorem{theorem}[definition]{Theorem}}{}
\makeatother
\def\P{\mathbb{P}}
\def\cM{{\mathcal{M}}}
\title[A Brill-Noether Theorem for (toric) surfaces]{A Brill-Noether Theorem for (toric) surfaces}
\author{Alessio Cela, Carl Lian}
\date{Source: arXiv:2503.18905v2 (CC BY 4.0)}
\begin{document}
\maketitle

\noindent\emph{Source and licence.} A Brill-Noether Theorem for (toric) surfaces. Version arXiv:2503.18905v2 (CC BY 4.0), distributed under the Creative Commons Attribution 4.0 International licence, as recorded in the arXiv licence metadata for this exact version and shown on the abstract page https://arxiv.org/abs/2503.18905v2. Full rights notice: licenses/arxiv-2503.18905-CC-BY-4.0.txt.\par\smallskip

\noindent\textit{Editorial scope.} Counted unit: the Proposition whose source label is \texttt{prop:multiple covers} in the article (source0.tex lines 384--388, source label \texttt{prop:multiple covers}), delivered together with the article's own complete original proof of it (source0.tex lines 390--418, one \texttt{proof} environment of 29 lines). No mathematics is written, repaired or re-derived: statement and proof are the article's own text line for line. Required context retained uncounted: severi\_thm (lines 356--366). Every definition, display and label of those lines is retained unchanged. Omitted from this package and not redistributed: 1--355, 367--383, 389--389, 419--711. Nothing inside a retained excerpt is omitted or altered: every retained body is delivered exactly as the article's lines are written, so no omission receipt of any kind accompanies this package. Citations: the retained excerpts carry no citation command at all, so no bibliography entry of the article is transcribed and no reference is redistributed. Reference adaptation: none was needed and none was made; every reference inside the delivered unit resolves inside it, so no unresolved reference or citation is produced. Printed slips kept literal: no printed word, symbol, hypothesis or number of the counted statement or of its proof was altered. Layout and class: the article's own class is \texttt{amsart} with packages including geometry, tikz, tikz-cd, url, hyperref, float, xy, amsmath, amssymb, enumitem, graphicx, mathdots, color, diagbox; the wrapper is the lane's pinned \texttt{amsart} editorial wrapper at 10pt on A4, which restates the theorem-like environments the retained text needs and adds the article's own macro definitions that the retained text needs (\texttt{\textbackslash P}, \texttt{\textbackslash cM}), with extra wrapper packages (bm, bbm, dsfont, xspace, cancel, mathtools). The delivered numbering is the unit's own. Assets: no retained span contains a figure environment, a graphics-inclusion command, a PDF-page inclusion, a table or a TikZ picture; no image file is copied into this package, the package declares no asset dependency and no image file is redistributed with it. Licence: version arXiv:2503.18905v2 of this article is distributed under the Creative Commons Attribution 4.0 International licence, as recorded in the arXiv licence metadata for this exact version and shown on the abstract page https://arxiv.org/abs/2503.18905v2. A full rights notice is delivered at licenses/arxiv-2503.18905-CC-BY-4.0.txt. Characters: every retained excerpt is pure ASCII, so the delivered engine, XeLaTeX with its default Latin Modern text font, reports no missing character.
\begin{theorem}\label{severi_thm}
    We have the following.
    \begin{enumerate}
     \item[(i)] If $\beta\cdot (-K_S)\ge 1$, then every component of $Y_{\beta,g}(S)$ has the expected dimension of 
     \begin{equation*}
         \beta\cdot (-K_S)+g-1.
     \end{equation*}
     \item[(ii)] If $\beta\cdot (-K_S)\ge 4$, then any general point of $Y_{\beta,g}(S)$ corresponds to a nodal curve $C\subset S$.
     \item[(iii)] If $C'\subset S$ is nodal and $\beta\cdot (-K_S)\ge 1$, then $Y_{\beta,g}(S)$ is smooth of expected dimension at $[C']$.
     \end{enumerate}
\end{theorem}

 \begin{proposition}\label{prop:multiple covers}
 Let $C$ be a general curve and let $S$ be a smooth, projective surface. Let $\beta\in H_2(S,\mathbb{Z})$ be non-zero, effective curve class. Let $\cM\subset \cM_\beta(C,S)$ be an irreducible component.

 Let $f:C\to S$ be a general point of $\cM$, and let $C'\subset S$ be the scheme-theoretic image of $f$. Suppose that $[C']\cdot (-K_S) \ge 4$. Then, $\cM$ is generically smooth of expected dimension.
 \end{proposition}

 \begin{proof}
Let $m$ be the degree of induced cover $g:C\to C'\subset S$. If $m=1$, then we are done by the above discussion, as $\cM_{\beta}(C,S)$ is isomorphic near $f$ to a general fiber of $\mu:Y_{\beta,g}(S)\to\cM_g$. We assume henceforth that $m>1$. In this case, the general curve $C$ admits no covers of a curve of positive geometric genus of degree $m>1$, unless $g=g(C')=1$. We consider the cases $g(C')=0$ and $g=g(C')=1$ separately.

%     \begin{enumerate}[label=\underline{Case \arabic*}]
%    \item  
Suppose first that $g(C')=0$. The moduli space of maps $f:C\to \P^1$ of degree $m$ has dimension $\rho(g,1,m)+\dim(\text{PGL}_2)=2m-g+1$. In particular, we need $\rho(g,1,m)=2m-2-g\ge0$.
        % \begin{equation}\label{eqn: estimate 1 in exc dim}
        % g\le 2m-2.
        % \end{equation}
By Theorem \ref{severi_thm}, the moduli space of embedded curves $C'\subset S$ in class $\beta':=[C']$ has dimension $\beta'\cdot (-K_S)-1$. Therefore, the component $\cM$ has dimension
        \begin{equation*}
            (2m-g+1)+(\beta'\cdot (-K_S)-1) \ge \beta\cdot (-K_S)-2g+2,
        \end{equation*}
        where the right hand side is the expected dimension. Rearranging and substituting $\beta=m\beta'$ gives
       \begin{equation*}
              (m-1)(\beta'\cdot (-K_S)-2) \le g.
         \end{equation*} 
         
Recall that we also have $g\le 2m-2$. Thus, the given multiple cover $f$ can only give rise to a component $\cM$ if $\beta'\cdot (-K_S)=4$ and $g=2m-2$, hence $\rho(g,1,m)=0$. In this case, the space $\cM_\beta(C,S)$ is isomorphic near $f$ to the 6-dimensional space of maps which factor as 
        \begin{equation*}
            C\to \P^1\to X,
        \end{equation*}
        where $f:C\to\P^1$ corresponds to one of a \emph{reduced} (for example, by the Gieseker-Petri theorem) set of points parametrizing rank 1 linear series on $C$ of minimal degree $m=\frac{g+2}{2}$, and the image of $\P^1\to S$ is a point of the Severi variety $Y_{\beta',0}(S)$. By Theorem \ref{severi_thm}, the Severi variety is generically smooth 
        of expected dimension, so $\cM$ is as well.
        
Consider now the exceptional case $g(C')=g(C)=1$. Since $C$ admits only finitely many maps of degree $m>1$ to some genus $1$ curve, by Theorem \ref{severi_thm}, the component $\cM$ would have dimension $\beta' \cdot (-K_S)$, which is less than the expected dimension $\beta\cdot (-K_S)$, a contradiction.

%\end{enumerate}
 \end{proof}

\end{document}
